% Abstract and the at-a-glance panel.
\section*{Summary}
\addcontentsline{toc}{section}{Summary}

Problem statement SIH26119 asks for an optimisation solver that India owns: one that reads the
models refinery planners already write, solves linear, mixed-integer and quadratic programmes,
uses GPUs where they help, and is not built on Xpress, CPLEX or any other existing solver library.
\nirnay{} is our answer. Every algorithm and every factorisation in it is implemented in the
repository: an MPS/QPS reader, presolve with exact primal and dual postsolve, scaling, a
Gilbert--Peierls sparse LU with eta updates and basis repair, an up-looking sparse Cholesky, a
quasi-definite $LDL\T$, a minimum-degree ordering, a bounded dual simplex, a Mehrotra
interior-point method, a restarted primal--dual hybrid gradient method (PDLP) on CPU and GPU, a
branch-and-bound code with reliability branching, propagation, heuristics and Gomory cuts, and a
QP interior-point method. HiGHS and Clarabel appear only in the benchmark harness, as independent
references that check every answer.

This report describes each algorithm as the code implements it, with its equations and
pseudocode; the numerical failures we hit and how each fix works; how answers are verified; and
the benchmark results, failures included. \Cref{sec:results} states plainly where \nirnay{} is behind:
it is slower than HiGHS, and it proves fewer MIPLIB instances optimal.

\begin{keybox}[At a glance (generated from \file{results/})]
\centering
\stat{\nSpx/\nSpxTot}{Netlib LPs solved by the dual simplex (objective within $10^{-6}$ of HiGHS)}\hfill
\stat{\nIpm/\nIpmTot}{Netlib LPs solved by the interior-point method}\hfill
\stat{\nQp/\nQpTot}{Maros--Mészáros convex QPs solved (within $10^{-6}$ of Clarabel)}\hfill
\stat{\nMip/\nMipTot}{MIPLIB\,3 instances proven optimal in \tlMip\,s\ifnum\nMipTot<\nMipList{} (sweep in progress)\fi}
\par\medskip
\stat{\ratioSpx$\times$}{dual simplex time vs.\ HiGHS, shifted geometric mean (shift \shiftA\,s)}\hfill
\stat{\gmrSpx$\times$}{dual simplex vs.\ HiGHS, geometric mean of per-instance ratios}\hfill
\stat{\largeSpeedMax$\times$}{largest GPU-over-CPU speed-up of our PDLP on the large LPs\ifnum\nLarge<\nLargeList{} run so far\fi}\hfill
\stat{\nDualBad}{dual-feasibility failures in \nDualOk{} Netlib solutions checked after postsolve}
\end{keybox}

\begin{honestbox}[Where we stand]
\nirnay{} is correct on the standard LP and QP test sets and substantially slower than a mature
open-source solver: the dual simplex takes \ratioSpx{} times as long as HiGHS in shifted geometric
mean over the \nSpxTot{} Netlib problems (\ratioSpxTenShift{} times with a \shiftB\,s shift, and
\gmrSpx{} times as the geometric mean of per-instance ratios). On MIPLIB\,3 with a \tlMip\,s limit,
branch-and-bound proves \nMip{} of the \nMipTot{} instances\ifnum\nMipTot<\nMipList{} run so far\fi{} optimal, against
\nMipRefOpt{} for HiGHS under the same limit. The engines are compiled with Numba inside the
inner loops, but the iteration drivers are Python, and the MIP code lacks most of the presolve,
cut families and heuristics that close gaps in production solvers.
\end{honestbox}
